% Options for packages loaded elsewhere
\PassOptionsToPackage{unicode}{hyperref}
\PassOptionsToPackage{hyphens}{url}
\PassOptionsToPackage{dvipsnames,svgnames,x11names}{xcolor}
\documentclass[
  10pt,
  a4paper,
]{article}
\usepackage{xcolor}
\usepackage[margin=18mm]{geometry}
\usepackage{amsmath,amssymb}
\setcounter{secnumdepth}{-\maxdimen} % remove section numbering
\usepackage{iftex}
\ifPDFTeX
  \usepackage[T1]{fontenc}
  \usepackage[utf8]{inputenc}
  \usepackage{textcomp} % provide euro and other symbols
\else % if luatex or xetex
  \usepackage{unicode-math} % this also loads fontspec
  \defaultfontfeatures{Scale=MatchLowercase}
  \defaultfontfeatures[\rmfamily]{Ligatures=TeX,Scale=1}
\fi
\usepackage{lmodern}
\ifPDFTeX\else
  % xetex/luatex font selection
  \setmainfont[]{DejaVu Serif}
  \setsansfont[]{DejaVu Sans}
  \setmonofont[]{DejaVu Sans Mono}
\fi
% Use upquote if available, for straight quotes in verbatim environments
\IfFileExists{upquote.sty}{\usepackage{upquote}}{}
\IfFileExists{microtype.sty}{% use microtype if available
  \usepackage[]{microtype}
  \UseMicrotypeSet[protrusion]{basicmath} % disable protrusion for tt fonts
}{}
\makeatletter
\@ifundefined{KOMAClassName}{% if non-KOMA class
  \IfFileExists{parskip.sty}{%
    \usepackage{parskip}
  }{% else
    \setlength{\parindent}{0pt}
    \setlength{\parskip}{6pt plus 2pt minus 1pt}}
}{% if KOMA class
  \KOMAoptions{parskip=half}}
\makeatother
\usepackage{longtable,booktabs,array}
\newcounter{none} % for unnumbered tables
\usepackage{calc} % for calculating minipage widths
% Correct order of tables after \paragraph or \subparagraph
\usepackage{etoolbox}
\makeatletter
\patchcmd\longtable{\par}{\if@noskipsec\mbox{}\fi\par}{}{}
\makeatother
% Allow footnotes in longtable head/foot
\IfFileExists{footnotehyper.sty}{\usepackage{footnotehyper}}{\usepackage{footnote}}
\makesavenoteenv{longtable}
\setlength{\emergencystretch}{3em} % prevent overfull lines
\providecommand{\tightlist}{%
  \setlength{\itemsep}{0pt}\setlength{\parskip}{0pt}}
\usepackage{xeCJK}
\setCJKmainfont[AutoFakeBold=1.4,AutoFakeSlant=0.2]{Droid Sans Fallback}
\setCJKsansfont[AutoFakeBold=1.4,AutoFakeSlant=0.2]{Droid Sans Fallback}
\setCJKmonofont{Droid Sans Fallback}
\emergencystretch=3em
\setlength{\parindent}{0pt}
\xeCJKDeclareCharClass{Default}{"2013,"00B7}
\usepackage{newunicodechar}
\newunicodechar{ℱ}{\ensuremath{\mathcal{F}}}
\newunicodechar{𝓕}{\ensuremath{\mathcal{F}}}
\newunicodechar{ℓ}{\ensuremath{\ell}}
\newunicodechar{≻}{\ensuremath{\succ}}
\renewcommand{\contentsname}{目录}
\usepackage{xurl}
\usepackage{graphicx}
\usepackage{pdflscape}
\AtBeginDocument{\hypersetup{colorlinks=true,linkcolor=black,urlcolor=blue,filecolor=blue,citecolor=black}}
\usepackage{bookmark}
\IfFileExists{xurl.sty}{\usepackage{xurl}}{} % add URL line breaks if available
\urlstyle{same}
\hypersetup{
  pdftitle={核心主线未完成项与定量复现执行计划},
  colorlinks=true,
  linkcolor={Maroon},
  filecolor={Maroon},
  citecolor={Blue},
  urlcolor={Blue},
  pdfcreator={LaTeX via pandoc}}

\title{核心主线未完成项与定量复现执行计划}
\author{}
\date{2026-09-12}

\begin{document}
\maketitle

本台账回答``还有哪几条核心claim没有完成，以及怎样继续完成严谨的底层原型机复现''。它区分已计算但与论文不符、仅在历史分支完成、缺少作者数值约定，以及原论文根本没有要求的更强定理。证据主文为\href{../MAJOR_CLAIMS_ZH.md}{本轮严格审计}，机器核验见\href{../FINAL_VERIFICATION.json}{FINAL\_VERIFICATION}。

\textbf{当前不能称论文稳健定量核心已经复现成功。}
已完成的是明确有限模型的求解、原式证书、同幅度IR→UV→共振信号、全部五组M/L交付和若干严格差异证明。最核心的未通过项仍是三代表的ρ稳健性及S0定量形状；其它未完成项必须按下列范围逐一理解。

\begingroup
\small
\setlength{\parskip}{0pt}
\setcounter{tocdepth}{2}
\tableofcontents
\endgroup
\clearpage

\section{一、六项尚未完整闭合的核心工作}\label{ux4e00ux516dux9879ux5c1aux672aux5b8cux6574ux95edux5408ux7684ux6838ux5fc3ux5de5ux4f5c}

\subsection{1.
三代表共同、合理的ρ：major，未通过}\label{ux4e09ux4ee3ux8868ux5171ux540cux5408ux7406ux7684ux3c1majorux672aux901aux8fc7}

已完成：三份完整中心的原式primal/支撑及进一步收紧；同一C/ImF/R的相位、η、散射强度、FF与电流谱；严格原生离散峰。新90°读数为803.391/701.034/696.485
MeV，最小ηP1约.777/.435/.242。精度提高未消除分散。

缺少：原文所报告三点在共同合理ρ能标附近的相近形状，及其近弹性解释的数值支持。精确xref上侧区域已证Im
S1(.7317
GeV)\textless−.50552051192，与原图该节点约30°的正虚部趋势不兼容；小坐标误差余量也不能消除符号差异。

判定边界：这是当前D及原几何区域的严格障碍，不是整个D内所有振幅、所有原参数实现或所有连续ρ极点的排除。下一步不能再只收紧相同问题的μ；须恢复或澄清会改变D/代表区域的真实来源细节。

\subsection{2. S0/S2的定量合理性及IR形状保留：S0
major；S2仅有比较差异}\label{s0s2ux7684ux5b9aux91cfux5408ux7406ux6027ux53cairux5f62ux72b6ux4fddux7559s0-majors2ux4ec5ux6709ux6bd4ux8f83ux5deeux5f02}

已完成：三代表S0/S2全C求值与论文对照，L/M比较；S0近tip全域下界Re
S0(1.1994
GeV)\textgreater.88103。该区域内45°--135°模180°的相位被排除，不是unwrap问题。

缺少：UV加入后仍与原文S0曲线及其低能主导解释相符。S2的已有曲线、有限样本和分辨率检查不能被误说成尚未计算，但部分定量差异仍保留；本轮未证明S2存在类似全域不相容。

下一步：与ρ使用同一份经来源核定的完整振幅，一起验收三波。不得分别从不同解挑出``好看的''S0与P1。

\subsection{3.
Fig.11的稳健性：五组计算完成，L比较未通过}\label{fig.11ux7684ux7a33ux5065ux6027ux4e94ux7ec4ux8ba1ux7b97ux5b8cux6210lux6bd4ux8f83ux672aux901aux8fc7}

已完成：M50/L8、L10、L12及M45/L10、M60/L10的全部要求配置。三个M50停止阈值已收紧到1e−6，均有中心与原式支撑。L跨度仍46.210
MeV，原图约12.539 MeV。

缺少：与原文相当的固定M的L稳定性；本方S0跨配置相对稳定，仍不等于其形状与论文一致。固定L的M跨度约54.064
MeV与原图54.349 MeV同量级，因此不能笼统称``M收敛全失败''。

下一步：先解决名义M50/L10的来源/代表问题，再按同一规则重做需要更新的配置。没有理由先扩大为大量M/L扫描，也不把数学M→∞收敛定理加入原文有限Appendix验收。

\subsection{4.
Fig.8的明显不对称及完整区域：定性收缩已通过，量化尚未对齐}\label{fig.8ux7684ux660eux663eux4e0dux5bf9ux79f0ux53caux5b8cux6574ux533aux57dfux5b9aux6027ux6536ux7f29ux5df2ux901aux8fc7ux91cfux5316ux5c1aux672aux5bf9ux9f50}

已完成：同xref的两侧严格收缩、上侧稍大；上/下比为{[}1.00493,1.09381{]}。冻结IR点的UV扩展已由支撑平面严格排除。

缺少：原图``上侧明显改变、下侧较少改变''的量化趋势；全部相关UV轮廓尚不能由这一个截面或五个点的凸包替代。原文没有给``明显''的统一误差阈值，不能把印刷图当作严格外排除证书。

下一步：在确定数值合同后重用同模型方向/内点，仅补最影响该比较的上下侧支撑；需要全域交付时再按已有几何误差证据补缺口。不能先扫描边界再按相移挑点。

\subsection{5.
Fig.3--6的原文对齐与同合同衔接：已有大量证据，不能一概说没做}\label{fig.36ux7684ux539fux6587ux5bf9ux9f50ux4e0eux540cux5408ux540cux8854ux63a5ux5df2ux6709ux5927ux91cfux8bc1ux636eux4e0dux80fdux4e00ux6982ux8bf4ux6ca1ux505a}

只读复审发现：历史\href{../../stage_B_delivery_20260907/regions/regions.json}{separate/PV区域}中pure有38份支持，几何误差约.00967577；六个ε区域均记录geometry\_ready=true，误差约.00907--.00985。这是存档所声明有限绘图域/metric下的完成状态，不因本轮没有重算便改称未完成。个别方向的宽对偶界也不自动推翻由其它平面封闭的整体几何界。

另一方面，后来的\href{../../mainline_alignment_20260909/B_combined/FIG4_CORE_CLAIMS_ZH.md}{combined历史支线}六域几何并未全部达标，且存在明显的原图范围差异。这与当前separate合同不同，不能混为一份状态。原图最大显示值不是作者的数学排除界。

\href{../../mainline_alignment_20260909/C_RESULT_ZH.md}{历史C1--C3}已交三色阈下行为、同C相移和比较，但其combined范数识别参考过Fig.5输出；不能追认证明本轮separate合同下的独立三色链。当前新的C3
IR对照已完成。缺少的是在最终来源合同下统一重接Fig.5→6→7及后续UV；不是缺少画图函数。蓝色全区间无零点等强命题尚未证明，但不能凭采样没有变号代替证明。

\subsection{6.
作者数值合同及整个原型身份：source-critical，尚未唯一恢复}\label{ux4f5cux8005ux6570ux503cux5408ux540cux53caux6574ux4e2aux539fux578bux8eabux4efdsource-criticalux5c1aux672aux552fux4e00ux6062ux590d}

已核对：原式变量/对称性、PV核、归一化、FESR量纲与截止、全联合变量、原生求值和原式证书。当前23模块的具名有限程序已自洽运行，不是低维替代模型。

尚未唯一恢复：2309作者的完整散射producer、χ范数/打包与可能权重、密度正则化约定，以及完整代表振幅选择。raw逐矩盒是原式直接读法；不能因后续程序改为三矩L2便偷偷替换。B/2的原式上界与原B见证严格分离；combined球严格排除当前tip
C，同时另有combined完整见证。这些说明约定有实际内容，但没有证明任一替代必然恢复论文相移。

数学上，仅说``某个范数''而不固定其归一化也不能唯一化容差：αN仍是范数，而\(\alpha N(v)\le\epsilon\)等价于\(N(v)\le\epsilon/\alpha\)。这说明需要补明什么，不是对论文数值结果的反证。

作者最初2403正文称额外幺正化在2309的1.2
GeV前作中不必要。故不能把后续幺正化当作原文已证遗漏的必做步骤。原文数值身份问题需要原始矩阵/参数/完整振幅或明确说明来解决；不能用拟合曲线补足。

\section{二、A1--F3逐项台账}\label{ux4e8ca1f3ux9010ux9879ux53f0ux8d26}

{\def\LTcaptype{none} % do not increment counter
\begin{longtable}[]{@{}
  >{\raggedright\arraybackslash}p{(\linewidth - 4\tabcolsep) * \real{0.2000}}
  >{\raggedright\arraybackslash}p{(\linewidth - 4\tabcolsep) * \real{0.3500}}
  >{\raggedright\arraybackslash}p{(\linewidth - 4\tabcolsep) * \real{0.4500}}@{}}
\toprule\noalign{}
\begin{minipage}[b]{\linewidth}\raggedright
步骤
\end{minipage} & \begin{minipage}[b]{\linewidth}\raggedright
当前完成度
\end{minipage} & \begin{minipage}[b]{\linewidth}\raggedright
精确未完成范围
\end{minipage} \\
\midrule\noalign{}
\endhead
\bottomrule\noalign{}
\endlastfoot
A1 交叉/色散/变量 & 明确有限模型内通过 &
作者完整producer身份；共同连续解析完成未声明 \\
A2 IR/UV输入 & 单位、显示公式及声明合同通过 &
作者χ/B/误差打包的唯一恢复；OPE/GMOR近似的理论余项没有严格全界 \\
A3 投影/求值 & 同H/C及原生精确源行核验通过 &
不移作另一插值或全能区证书 \\
B1 支撑/求解器 & 本轮更严格支撑、中心和下降修补通过 &
精确驻点区间证明未要求/未提供；数值中心范围明确 \\
B2 Fig.3 & 历史有限几何已有通过记录 &
原作者合同及原图定量身份尚未闭合 \\
B3 Fig.4 & 历史sep六域有几何通过；combined另记 &
最终合同的原图范围/全部比较；不能混用两分支 \\
C1 Fig.5 & 历史三色定性/零点证据保留 &
最终合同下独立三色链；蓝色无零点强命题不得伪造 \\
C2 Fig.6 & 既有完整幅度图已交 & 与最终C1/C3/UV同合同重新衔接 \\
C3 Fig.7 & 新完整IR对照通过 & 无ρ是对照结果，不是待修绘图缺陷 \\
D1 Gram/FF & 全部联合变量及归一化通过 &
不等于两π谱饱和、近弹性或全谱完成 \\
D2 FESR/高能FF & 当前四raw盒及FF条件通过 &
作者未明示口径；近似QCD输入不等于精确QCD定理 \\
D3 联合可行性 & M50三点与五组均完成 & 不等于所选幅度具有正确定量ρ \\
E1 Fig.8 & 精确截面弱不对称通过 & 明显不对称及完整UV轮廓对齐 \\
E2 代表选择 & 独立冻结三点通过 & 作者最终选择/tie-break尚未恢复 \\
E3 Fig.9--10 & 所有曲线/谱/峰已计算 &
ρ稳健性和S0未通过；两项有近区域全域障碍 \\
F1 五配置 & 全部完成 & 若合同改变，必须相应重验 \\
F2 稳定性 & 已量化并收紧M50精度 & L稳健性未对齐；M不能笼统判失败 \\
F3 交付 & 来源、失败、证明、数据、图齐备 &
交付齐备不等于全Fig.3--11物理复现成功 \\
\end{longtable}
}

\section{三、下一阶段顺序计划}\label{ux4e09ux4e0bux4e00ux9636ux6bb5ux987aux5e8fux8ba1ux5212}

\subsection{P0：锁定验收与来源，先于任何新优化}\label{p0ux9501ux5b9aux9a8cux6536ux4e0eux6765ux6e90ux5148ux4e8eux4efbux4f55ux65b0ux4f18ux5316}

保留本轮D及全部失败为基线。建立一份逐条输入表：原文显式给定、由公式严格换算、后续有依据但2309未确认、完全未明示。所需补明项仅限上文χ/B/误差/散射producer/选择，不展开其它现象学输入。

验收分两层：数值层要求全变量原式primal/支撑、同幅度身份和声明过的数值中心；物理层比较三点ρ能标/峰形、η、S0/S2、Fig.8及原M/L配置。论文未给统一误差阈值的项目先声明比较指标与来源精度，不在结果出来后为通过而改阈值，也不把每个印刷点精确重合当要求。

\textbf{完成条件：}
每个会改变D或代表区域的选择都有可核验来源或明确的条件模型身份。若2309资料仍缺失，应明确``无法唯一复原作者程序''，不能宣称已经完成这一关。

\subsection{P1：最小算子/完整幅度核对}\label{p1ux6700ux5c0fux7b97ux5b50ux5b8cux6574ux5e45ux5ea6ux6838ux5bf9}

若取得作者2309矩阵或完整点，先重验而不优化：确认变量顺序、ρ2非对角系数、减除坐标、κ/FF归一化、FESR误差集合与目标、节点/波数及投影泛函。先用常数/单密度/双密度的独立代数检查定位，再检查相关完整行或矩阵。

需要MMA时可用已提供容器，以WL实现独立公式核对，并将输出数据/来源接入同一CLI。后续2403程序只能作明确标记的公式核对；改ν0、通道或误差集合后的矩阵不能冒充2309原producer。若现有算子已经满足所需核对，不重复建立新平台。

\textbf{完成条件：}
有差异便给出明确代数来源及等价/非等价证明；无差异则排除此实现解释。作者点若违反当前原式，应列出具体裕量/单位/支持位置，不以其曲线正确为理由放宽约束。

\subsection{P2：有来源依据后修补模型并完成名义M50/L10}\label{p2ux6709ux6765ux6e90ux4f9dux636eux540eux4feeux8865ux6a21ux578bux5e76ux5b8cux6210ux540dux4e49m50l10}

只实施P1已定位的真实差异。等价坐标变换必须同时搬运完整约束；非等价更改必须声明新D。复用旧完整点前重新primal审计，不继承旧支持或中心；必要时Phase-I，随后完成新的+x和固定截面支持。

代表规则先于相位冻结。若原规则仍是本轮D内的相同近tip/精确xref上侧区域，则本轮两个负上界仍有效，不能再期待同样优化得到相反相位。若新来源合法改变D或区域，要明确指出旧证书为何不再适用。

\textbf{完成条件：}
全部C/ImF/R、所有原生盘/Gram/χ/L4/FESR/FF通过；中心身份可回查，支持间隙达标；无修改后曲线反选输入。

\subsection{P3：同一幅度完成C→D→E的物理验收}\label{p3ux540cux4e00ux5e45ux5ea6ux5b8cux6210cdeux7684ux7269ux7406ux9a8cux6536}

重新衔接最终合同下的三色阈下图、同C物理IR对照及三UV代表。对每份完整点报告δ、η、强度、FF与两π/总电流谱，保留所有局部峰。采用原文90°识别的数值范围，同时明确它不是唯一极点证明。

重新运行本轮两个预定分量探针，作为很敏感的审计关卡。若一份被称为相同D/相同区域的新结果显示相反符号，先检查算子、目标、C身份及证书，不宣布物理突破。S0与P1必须来自同一完整解，不能拼接。

\textbf{完成条件：}
原论文稳健定量趋势得到实际数据支持，并在报告中给出剩余偏差/近似来源。只出现单个峰或只有tip较好不满足三点稳健性要求。

\subsection{P4：按最终合同交付区域和五组稳定性}\label{p4ux6309ux6700ux7ec8ux5408ux540cux4ea4ux4ed8ux533aux57dfux548cux4e94ux7ec4ux7a33ux5b9aux6027}

同D时优先复用已认证几何/支撑；D改变时先重验。Fig.8先补决定收缩比较的最小方向，再按几何误差定位剩余轮廓缺口。Fig.3/4同样区分已有geometry\_ready与原图身份，不重复扫描已经完成的区域。

使用同一预先固定的代表规则完成原五组M/L；M变化同时改变变量、节点和FESR权重，不能转移旧可行性。原M50名义点尚未物理对齐时，不先扩大M/L扫参。

\textbf{完成条件：}
有完整有限区域/比较范围说明；三波与ρ指标的M/L变化支持原文相应稳健性；未通过则继续保留差异，不用配置挑选遮盖。

\subsection{P5：最后的独立审计与停止规则}\label{p5ux6700ux540eux7684ux72ecux7acbux5ba1ux8ba1ux4e0eux505cux6b62ux89c4ux5219}

检查所有图绑定同一完整幅度、所有输入和producer可追溯、所有原式上下界与失败保留。测试限于独立数学必要检查，继续保持少量职责模块和最多三个测试文件。科学计算逐个、每次有运行上限，只在结果给出具体进展时续算。

只有上述来源、数值和物理层同时满足，才报告稳健定量核心复现成功。若来源未唯一恢复，结论必须保持条件性；若在来源明确的同一问题上仍有严格不相容，应提交条件充分的复核证据，而不是继续更换误差范数、B、质量、截止或饱和条件来贴图。

\section{四、暂不加入主线的更强目标}\label{ux56dbux6682ux4e0dux52a0ux5165ux4e3bux7ebfux7684ux66f4ux5f3aux76eeux6807}

唯一第二张面极点、全能量/无限自旋认证、整个可行集都必有ρ、精确驻点区间Newton证明及B→∞平台，均不自动成为2309有限原型的新增前提。局部矩阵投影非闭、模不唯一决定相位等数学问题单列保留；它们不能冒充整篇论文的反例。

\textbf{下一项实际物理决策是P0/P1中的数值合同唯一化与有针对性的算子/完整点核对。}
当前模型在指定区域的符号障碍已经严格成立，继续相同μ迭代不能代替这一步；未获得新来源时，不承诺可以在完全不变的合同中得到已被排除的曲线。

\end{document}
